\section{Method}
CertiPatch uses a gated low-rank hookpoint patch (GLR-HP) on the residual stream at the answer position.
For candidate layers $\ell\in\mathcal{L}$:
\[
h_{\ell,p} \leftarrow h_{\ell,p} + s(x)\,U_\ell(V_\ell^\top h_{\ell,p}),
\]
with rank $r$ and candidate layers determined by configuration.

\paragraph{Constrained minimality.}
We optimize:
\[
\min_{\Delta\phi}\; \mathcal{L}_{\text{collateral}}(\Delta\phi) + \lambda_{\text{reg}}\mathcal{R}(\Delta\phi)
\quad \text{s.t.}\quad g_{\text{spec}}(\Delta\phi)=0,
\]
where $g_{\text{spec}}$ encodes in-scope spec violations.
The inner solver uses an augmented Lagrangian schedule.
The outer loop is counterexample-guided and adds hardest violations to the active set each round.
This CEGIS-style structure is used for closure under discovered failures while keeping the patch sparse and localized.

\paragraph{Collateral objectives.}
Collateral is measured on dedicated reference suites: short-form Yes/No KL and long-horizon generation drift.
These suites are gate-aware: in-scope prompts are evaluated as such, and out-of-scope prompts remain excluded by construction.
The resulting optimization target is explicitly ``repair under collateral budget'', not unconstrained task adaptation.

\paragraph{Compositionality.}
We train separate patches for two specs and evaluate interference under additive composition and sequential repair (A$\to$B and B$\to$A).
